% GENERATED FILE. Edit canonical lessons, then run npm run generate:books.
% canonical-source-hash: fnv1a64:eb7e4a101a236e3e
% canonical-lessons: RU-C377-perenosit, RU-C377-peredumyvat, RU-C377-otkazyvatsya, RU-C377-naznachat, RU-C377-sobiratsya

\chapter{Changing plans}
\label{ch:ru-changing-plans}
\hlchaptermodality{377}

\begin{chapteropening}
\textbf{By the end of this chapter:} \emph{I can move a meeting to another day, change my mind, turn something down, fix a time and say what I am going to do.}

Complete the last lesson of chapter 377: I can move a meeting to another day, change my mind, turn something down, fix a time and say what I am going to do.
\end{chapteropening}

\section[\texorpdfstring{\ru{переносить} (perenosít)}{переносить (perenosít)}]{\ru{переносить} (perenosít) — to postpone, to move (to another time); to carry over}

\label{lesson:RU-C377-perenosit}



\begin{quote}
Before the new one: say the Russian for at least, then the Russian for in case.
\end{quote}

\subsection*{You'll want to know: \ru{переносить}}
\textbf{\ru{переносить}} (\emph{perenosít}) — \textquotedblleft{}to postpone, to move (to another time); to carry over\textquotedblright{}.

\begin{quote}
\ru{Мы} \ru{переносим} \ru{встречу} \ru{на} \ru{понедельник}.
\end{quote}

\begin{quote}
\emph{My perenósim vstréchu na ponedélnik.}
\end{quote}

\begin{quote}
We are moving the meeting to Monday.
\end{quote}

Imperfective; present \emph{\ru{переношу}, \ru{переносишь}}; perfective \emph{\ru{перенести}}. \emph{\ru{Переносить} \ru{на}} + accusative — to move to. It also means to bear: \emph{\ru{плохо} \ru{переносить} \ru{холод}}.

\begin{cousinweb}
\textbf{\ru{переносить}} is built from \emph{\ru{пере}-} \textquotedblleft{}across, over\textquotedblright{} and \emph{\ru{носить}} \textquotedblleft{}to carry\textquotedblright{}, inherited from Proto-Slavic: to carry a meeting over to another day.
\end{cousinweb}

\begin{grammarlens}[title={What you've built}]
One more word for telling, explaining and describing at B1.
\end{grammarlens}

\subsection*{Your turn}
\begin{itemize}
\raggedright
  \item \emph{Say it:} \emph{perenosít}
  \item \emph{Say it:} \emph{perenosít}, once more
  \item \emph{Say it:} a sentence of your own with \emph{perenosít}
  \item \emph{Recall:} say the Russian for at least, then the Russian for in case, then say \emph{perenosít} again
\end{itemize}

\begin{tcolorbox}[breakable,colback=teal!4,colframe=teal!35!black,title={Before you move on}]
What does \ru{переносить} mean? (\textquotedblleft{}To postpone, to move (to another time); to carry over\textquotedblright{}.) Say it once more.
\end{tcolorbox}
\section[\texorpdfstring{\ru{передумывать} (peredúmyvat)}{передумывать (peredúmyvat)}]{\ru{передумывать} (peredúmyvat) — to change one’s mind}

\label{lesson:RU-C377-peredumyvat}



\begin{quote}
Before the new one: say the Russian for in case, then the Russian for to postpone.
\end{quote}

\subsection*{You'll want to know: \ru{передумывать}}
\textbf{\ru{передумывать}} (\emph{peredúmyvat}) — \textquotedblleft{}to change one’s mind\textquotedblright{}.

\begin{quote}
\ru{Он} \ru{часто} \ru{передумывает} \ru{в} \ru{последнюю} \ru{минуту}.
\end{quote}

\begin{quote}
\emph{On chásto peredúmyvayet v poslédnyuyu minútu.}
\end{quote}

\begin{quote}
He often changes his mind at the last minute.
\end{quote}

Imperfective; perfective \emph{\ru{передумать}}: \emph{\ru{я} \ru{передумал}} — I have changed my mind.

\begin{cousinweb}
\textbf{\ru{передумывать}} is built from \emph{\ru{пере}-} \textquotedblleft{}over again\textquotedblright{} and \emph{\ru{думать}} \textquotedblleft{}to think\textquotedblright{}: to think a thing over again and come out the other way.
\end{cousinweb}

\begin{grammarlens}[title={What you've built}]
One more word for telling, explaining and describing at B1.
\end{grammarlens}

\subsection*{Your turn}
\begin{itemize}
\raggedright
  \item \emph{Say it:} \emph{peredúmyvat}
  \item \emph{Say it:} \emph{peredúmyvat}, once more
  \item \emph{Say it:} a sentence of your own with \emph{peredúmyvat}
  \item \emph{Recall:} say the Russian for in case, then the Russian for to postpone, then say \emph{peredúmyvat} again
\end{itemize}

\begin{tcolorbox}[breakable,colback=teal!4,colframe=teal!35!black,title={Before you move on}]
What does \ru{передумывать} mean? (\textquotedblleft{}To change one’s mind\textquotedblright{}.) Say it once more.
\end{tcolorbox}
\section[\texorpdfstring{\ru{отказываться} (otkázyvatsya)}{отказываться (otkázyvatsya)}]{\ru{отказываться} (otkázyvatsya) — to refuse, to turn down; to give up}

\label{lesson:RU-C377-otkazyvatsya}



\begin{quote}
Before the new one: say the Russian for to postpone, then the Russian for to change one’s mind.
\end{quote}

\subsection*{You'll want to know: \ru{отказываться}}
\textbf{\ru{отказываться}} (\emph{otkázyvatsya}) — \textquotedblleft{}to refuse, to turn down; to give up\textquotedblright{}.

\begin{quote}
\ru{Она} \ru{отказывается} \ru{от} \ru{сладкого}.
\end{quote}

\begin{quote}
\emph{Oná otkázyvayetsya ot sládkogo.}
\end{quote}

\begin{quote}
She is giving up sweet things.
\end{quote}

Imperfective; perfective \emph{\ru{отказаться}}. \emph{\ru{Отказываться} \ru{от}} + genitive (to turn down or give up), or + infinitive: \emph{\ru{отказался} \ru{помочь}}.

\begin{cousinweb}
\textbf{\ru{отказываться}} is built from \emph{\ru{от}-} \textquotedblleft{}away\textquotedblright{} and \emph{\ru{казать}} \textquotedblleft{}to show, to say\textquotedblright{}, Proto-Slavic \emph{*kazati} (EtymDB): to say a thing away from yourself.
\end{cousinweb}

\begin{grammarlens}[title={What you've built}]
One more word for telling, explaining and describing at B1.
\end{grammarlens}

\subsection*{Your turn}
\begin{itemize}
\raggedright
  \item \emph{Say it:} \emph{otkázyvatsya}
  \item \emph{Say it:} \emph{otkázyvatsya}, once more
  \item \emph{Say it:} a sentence of your own with \emph{otkázyvatsya}
  \item \emph{Recall:} say the Russian for to postpone, then the Russian for to change one’s mind, then say \emph{otkázyvatsya} again
\end{itemize}

\begin{tcolorbox}[breakable,colback=teal!4,colframe=teal!35!black,title={Before you move on}]
What does \ru{отказываться} mean? (\textquotedblleft{}To refuse, to turn down; to give up\textquotedblright{}.) Say it once more.
\end{tcolorbox}
\section[\texorpdfstring{\ru{назначать} (naznachát)}{назначать (naznachát)}]{\ru{назначать} (naznachát) — to set, to fix (a time); to appoint}

\label{lesson:RU-C377-naznachat}



\begin{quote}
Before the new one: say the Russian for to change one’s mind, then the Russian for to refuse.
\end{quote}

\subsection*{You'll want to know: \ru{назначать}}
\textbf{\ru{назначать}} (\emph{naznachát}) — \textquotedblleft{}to set, to fix (a time); to appoint\textquotedblright{}.

\begin{quote}
\ru{Мы} \ru{назначили} \ru{встречу} \ru{на} \ru{пять} \ru{часов}.
\end{quote}

\begin{quote}
\emph{My naznáchili vstréchu na pyat chasóv.}
\end{quote}

\begin{quote}
We set the meeting for five o’clock.
\end{quote}

Imperfective; perfective \emph{\ru{назначить}}. It also means to appoint someone to a post, or to prescribe a medicine.

\begin{cousinweb}
\textbf{\ru{назначать}} is built from \emph{\ru{на}-} and \emph{\ru{знак}} \textquotedblleft{}a sign, a mark\textquotedblright{}, on the root of \emph{\ru{знать}} \textquotedblleft{}to know\textquotedblright{}, Proto-Slavic \emph{*znati}: to mark a time out.
\end{cousinweb}

\begin{grammarlens}[title={What you've built}]
One more word for telling, explaining and describing at B1.
\end{grammarlens}

\subsection*{Your turn}
\begin{itemize}
\raggedright
  \item \emph{Say it:} \emph{naznachát}
  \item \emph{Say it:} \emph{naznachát}, once more
  \item \emph{Say it:} a sentence of your own with \emph{naznachát}
  \item \emph{Recall:} say the Russian for to change one’s mind, then the Russian for to refuse, then say \emph{naznachát} again
\end{itemize}

\begin{tcolorbox}[breakable,colback=teal!4,colframe=teal!35!black,title={Before you move on}]
What does \ru{назначать} mean? (\textquotedblleft{}To set, to fix (a time); to appoint\textquotedblright{}.) Say it once more.
\end{tcolorbox}
\section[\texorpdfstring{\ru{собираться} (sobirátsya)}{собираться (sobirátsya)}]{\ru{собираться} (sobirátsya) — to be going to, to intend; to gather}

\label{lesson:RU-C377-sobiratsya}



\begin{quote}
Before the new one: say the Russian for to refuse, then the Russian for to set.
\end{quote}

\subsection*{You'll want to know: \ru{собираться}}
\textbf{\ru{собираться}} (\emph{sobirátsya}) — \textquotedblleft{}to be going to, to intend; to gather\textquotedblright{}.

\begin{quote}
\ru{Летом} \ru{мы} \ru{собираемся} \ru{на} \ru{море}.
\end{quote}

\begin{quote}
\emph{Létom my sobiráyemsya na móre.}
\end{quote}

\begin{quote}
In the summer we are planning to go to the seaside.
\end{quote}

Imperfective; perfective \emph{\ru{собраться}}. With an infinitive, to intend: \emph{\ru{я} \ru{собираюсь} \ru{уехать}}. It also means to gather: \emph{\ru{гости} \ru{собираются}}.

\begin{cousinweb}
\textbf{\ru{собираться}} is the reflexive of \emph{\ru{собирать}} \textquotedblleft{}to gather\textquotedblright{}, from \emph{\ru{со}-} \textquotedblleft{}together\textquotedblright{} and \emph{\ru{бирать}}, from \emph{\ru{брать}} \textquotedblleft{}to take\textquotedblright{}: to gather yourself up for a journey or a plan.
\end{cousinweb}

\begin{grammarlens}[title={What you've built}]
One more word for telling, explaining and describing at B1.
\end{grammarlens}

\subsection*{Your turn}
\begin{itemize}
\raggedright
  \item \emph{Say it:} \emph{sobirátsya}
  \item \emph{Say it:} \emph{sobirátsya}, once more
  \item \emph{Say it:} a sentence of your own with \emph{sobirátsya}
  \item \emph{Recall:} say the Russian for to refuse, then the Russian for to set, then say \emph{sobirátsya} again
\end{itemize}

\begin{tcolorbox}[breakable,colback=teal!4,colframe=teal!35!black,title={Before you move on}]
What does \ru{собираться} mean? (\textquotedblleft{}To be going to, to intend; to gather\textquotedblright{}.) Say it once more.
\end{tcolorbox}
